\originalsection{18.315：作业1（原卷 Fall 2005）}
本组收录原题1、2、3、5, 共七个小问单元. 网格图 $G_{m,n}$ 的顶点为
$[m]\times[n]$, 两点相邻当且仅当在一个坐标上相差1、另一个坐标相同;
即两条分别有 $m,n$ 个顶点的路径的笛卡儿积. 以下颜色均有固定名称.

\begin{exercise}\label{mit:18315-hw1-q1}
\textbf{18.315, 原卷 Fall 2005, 作业1, 题1.}
设 $c(n)$ 为 $G_{n,n}$ 的正常三染色数.
\begin{enumerate}[label=(\alph*)]
\item 证明存在常数 $C,\varepsilon>0$, 使 $c(n)>C(1+\varepsilon)^{n^2}$.
\item 证明 $\log c(n)/n^2$ 收敛到某个正数 $\alpha$.
\end{enumerate}
\end{exercise}
\sourceof{官方 HW1 第1页题1; 以下为编者独立解答, 未有官方答案}
\begin{solution}
(a) 按两坐标之和的奇偶性二分顶点. 把较小的一部分全部染颜色1,
较大的一部分的每个顶点独立选择颜色2或3. 同一部分没有边, 因而每种选择都合法.
这样至少得到 $2^{\lceil n^2/2\rceil}$ 种染色. 可取 $C=1/2$ 与
$1+\varepsilon=\sqrt2$, 此时所需严格不等式对每个 $n\ge1$ 成立.

(b) 令 $a_n=\log c(n)/n^2$, $\alpha_* =\inf_{m\ge1}a_m$.
上一步及逐点选择三色的上界给出
$\tfrac12\log2\le\alpha_*\le a_n\le\log3$.
固定正整数 $m$, 令 $q=\lfloor n/m\rfloor$.
把大网格分成 $q^2$ 个互不重叠的 $m\times m$ 顶点块及剩余顶点.
一个正常染色限制在每个块上仍是正常染色; 忽略块间的边只会增加可选数.
剩余顶点可以各自最多选择三种颜色, 所以
\[
c(n)\le c(m)^{q^2}3^{n^2-q^2m^2}.
\]
取对数、除以 $n^2$, 再令 $n\to\infty$, 得 $\limsup_n a_n\le a_m$.
这对每个 $m$ 都成立, 故 $\limsup_n a_n\le\alpha_*$.
另一方面每个 $a_n\ge\alpha_*$, 所以 $a_n\to\alpha_* >0$.
这里不需要把任意小块染色拼成合法大块染色; 限制映射的上界已足够.
\end{solution}

\begin{exercise}\label{mit:18315-hw1-q2}
\textbf{18.315, 原卷 Fall 2005, 作业1, 题2.}
记 $N_k$ 为 $G_{n,n}$ 的正常 $k$ 染色数. 在相对误差不超过10\%的意义下近似计算:
(a) $n=100,k=1\,000\,000$; (b) $n=100,k=1000$.
\end{exercise}
\sourceof{官方 HW1 第1页题2; 编者独立解答, 误差界另行严格推导}
\begin{solution}
按行从左到右依次加入顶点. 第一个顶点有 $k$ 种选择;
第一行、第一列其余 $2(n-1)$ 个顶点各有 $k-1$ 种;
内部顶点前面已有左、上两个邻点, 至少有 $k-2$ 种. 因而
\[
L_k=k(k-1)^{198}(k-2)^{9801}\le N_k \qquad(n=100).
\]
还须控制两个邻点同色而多出来的选择, 不能仅作独立事件近似.
在某个内部顶点加入前, 对已有前缀图的全部正常染色均匀取样.
记 $u$ 为左邻点, $v$ 为上邻点. 在前缀图中 $u$ 至多有两个邻点,
且 $u,v$ 不相邻. 固定除 $u$ 外所有颜色后, $u$ 至少有 $k-2$ 种合法颜色;
因此条件概率 $\Prob(\chi(u)=\chi(v))\le1/(k-2)$.
平均后同一上界仍成立. 于是这个内部顶点的平均延伸数介于
$k-2$ 和 $k-2+1/(k-2)$ 之间. 对9801个内部步骤逐次相乘, 得
\[
1\le \frac{N_k}{L_k}
\le\left(1+\frac1{(k-2)^2}\right)^{9801}
\le\exp\!\left(\frac{9801}{(k-2)^2}\right).
\]
两问都可直接取 $L_k$ 为近似值. 对(a), 相对低估不足 $10^{-8}$;
对(b), 相对低估不足 $0.0099$, 均比10\%要求更强.
为便于读数, 其数量级分别为
\[
L_{1\,000\,000}\approx9.803947\times10^{59999},\qquad
L_{1000}\approx2.468322\times10^{29991}.
\]
精确可复核的近似定义是前面的整数幂乘积 $L_k$, 末行只用于展示大小.
\end{solution}

\begin{exercise}\label{mit:18315-hw1-q3}
\textbf{18.315, 原卷 Fall 2005, 作业1, 题3.}
设 $S_k(n)$ 为网格图 $G_{n,n}$ 的正常 $k$ 染色集.
证明任意两个染色可通过每次只改变一个顶点颜色、始终保持正常染色而相互转化:
(a) $k=5$; (b) $k=4$.
\end{exercise}
\sourceof{官方 HW1 第1页题3; 编者独立解答}
\begin{solution}
网格图是2退化的: 在任意非空诱导子图中, 取行坐标最小、再取该行列坐标最小的顶点,
它没有上方或左方邻点, 所以度数至多2.
下面直接证明任意2退化图在至少四种颜色时的重新配置连通性.

对顶点数归纳. 删去一个度数至多2的顶点 $x$,
归纳得到其余顶点从初始染色到目标染色的逐点改色序列.
模拟每一步. 若准备给 $x$ 的某个邻点赋予的颜色恰等于 $x$ 的现色,
先改 $x$: 避开它至多两个邻点的现色, 再避开即将使用的颜色,
总共最多禁止三色, 四种颜色中至少还剩一种.
先给 $x$ 赋这个颜色, 再执行邻点的改色, 两步都正常.
其余步骤无需改 $x$. 当删去 $x$ 的图达到目标后,
直接给 $x$ 赋目标颜色, 因目标染色本来正常, 最后一步也合法.
单顶点情形显然, 因而归纳完成. 同一构造同时解决 $k=4$ 和 $k=5$.
\end{solution}

\begin{exercise}\label{mit:18315-hw1-q5}
\textbf{18.315, 原卷 Fall 2005, 作业1, 题5.}
从有 $n$ 个给定顶点、恰有 $m=2n$ 条边的简单图中均匀随机选择 $H$.
它更可能是二部图, 还是更可能不是二部图?
\end{exercise}
\sourceof{官方 HW1 第1页题5; 编者明确均匀选边模型, 原题可行参数须 $n\ge5$}
\begin{solution}
当 $n<5$ 时 $2n>\binom n2$, 所述随机模型没有样本; 以下 $n\ge5$.
固定一个顶点二分, 跨边至多 $B=\lfloor n^2/4\rfloor$, 而全部候选边为 $M=\binom n2$.
所有 $2n$ 条边都跨越该二分的概率为 $\binom B{2n}/\binom M{2n}$,
若 $B<2n$ 则为零; 其余情况下逐因子比较得不超过 $(B/M)^{2n}$.
无序二分至多 $2^{n-1}$ 个. 任何二部图至少由一个二分见证, 故并集界给出
\[
\Prob(H\text{ 二部})\le
2^{n-1}\left(\frac{n}{2(n-1)}\right)^{2n}
=\frac12\left(\frac{n^2}{2(n-1)^2}\right)^n
\le\frac12\left(\frac{25}{32}\right)^n<\frac12.
\]
因此对每个可行的 $n$, 非二部图更可能; 随 $n$ 增大, 二部概率还呈指数下降上界.
\end{solution}
